\documentclass[UTF8]{ctexart}
\usepackage{xeCJK}
\usepackage{geometry}
\geometry{left=2cm,right=2cm,top=2.5cm,bottom=2.5cm}
\setlength{\headheight}{12.7pt}

\usepackage{titlesec}
\titleformat{\section}
  {\normalfont\large\bfseries}
  {\thesection}
  {1em}
  {}
\titlespacing*{\section}
  {0pt}
  {3.5ex plus 1ex minus .2ex}
  {2.3ex plus .2ex}

\usepackage{amsmath}
\usepackage{amssymb}
\usepackage{amsthm}
\usepackage{enumitem}
\setlist[enumerate]{itemsep=0pt,topsep=0pt,parsep=0pt,leftmargin=0pt,itemindent=2.5em}
\setlist[itemize]{itemsep=0pt,topsep=0pt,parsep=0pt,leftmargin=0pt,itemindent=2.5em}
\usepackage{fancyhdr}
\pagestyle{fancy}
\fancyhf{}
\usepackage{graphicx}
\usepackage{hyperref}
\usepackage{indentfirst}
\usepackage{lmodern}


\begin{document}
\setlength{\abovedisplayskip}{1pt}
\setlength{\belowdisplayskip}{1pt}

\section{子结构}

\hypertarget{sec:定义1.1}{}
\noindent\textbf{定义1.1 (子结构)}

给定一阶语言$\mathcal{L}$和
$\mathcal{L}$-\href{01 一阶语言的结构#nameddest=sec:定义1.1}{结构} $A,B$. 如果:
\begin{enumerate}
    \item $A\subset B$,
    \item 对任意的常元$c$, $c^A=c^B$,
    \item 对任意的$n$元函数符号$f$, $f^A=f^B\!\restriction_{A^n}$,
    \item 对任意的$n$元关系符号$r$, $r^A=r^B\cap A^n$;
\end{enumerate}
就称$A$是$B$的$\mathcal{L}$-\textbf{子结构}, 记作$A\leqslant B$.

\vspace{10pt}

\hypertarget{sec:定理1.1}{}
\noindent\textbf{定理1.1}

子结构关系是偏序关系.

\vspace{10pt}

显然子结构的所有解释由其父结构所确定. 下面的定理给出了什么子集可以成为子结构.

\vspace{10pt}

\hypertarget{sec:定理1.2}{}
\noindent\textbf{定理1.2}

给定一阶语言$\mathcal{L}$和$\mathcal{L}$-结构$B$.
对任意非空的$A\subset B$, $A\leqslant B$当且仅当:
\begin{enumerate}
    \item 对任意的常元$c$, $c^B\in A$,
    \item 对任意的$n$元函数符号$f$和$a_1,\cdots,a_n\in A$, $f^B(a_1,\cdots,a_n)\in A$.
\end{enumerate}

\noindent{\kaishu 证明.}

\noindent{\kaishu 必要性.}

假设$A\leqslant B$, 那么:
\begin{enumerate}
    \item 对任意的常元$c$, $c^B=c^A\in A$,
    \item 对任意的$n$元函数符号$f$和$a_1,\cdots,a_n\in A$,
    $f^B(a_1,\cdots,a_n)=f^A(a_1,\cdots,a_n)\in A$.
\end{enumerate}

\noindent{\kaishu 充分性.}

假设$A\subset B$含有所有的$c^B$且对所有的$f^B$封闭, 那么:
\begin{enumerate}
    \item 对任意的常元$c$, 定义$c^A=c^B$,
    \item 对任意的$n$元函数符号$f$, 定义$f^A=f^B\!\restriction_{A^n}$,
    \item 对任意的$n$元关系符号$r$, 定义$r^A=r^B\cap A^n$,
\end{enumerate}
显然就有$A\leqslant B$. \hfill $\qedsymbol$





\newpage

\section{子结构与同态}

\hypertarget{sec:定理2.1}{}
\noindent\textbf{定理2.1 (嵌入同态)}

给定一阶语言$\mathcal{L}$和$\mathcal{L}$-结构$A,B$.
如果$A\leqslant B$, 那么嵌入映射
\[
    \iota:A\to B,\,a\mapsto a
\]
是从$A$到$B$的单同态, 称之为\textbf{嵌入同态}.

\vspace{10pt}

一般同态的值域也是子结构.

\vspace{10pt}

\hypertarget{sec:定理2.2}{}
\noindent\textbf{定理2.2}

给定一阶语言$\mathcal{L}$和$\mathcal{L}$-结构$A,B$.

设$p:A\to B$是同态, 则存在唯一的$B'\subset B$和$q:A\to B'$使得:
\begin{enumerate}
    \item $B'\leqslant B$,
    \item $q$是从$A$到$B'$的满同态,
    \item $p=\iota\circ q$.
\end{enumerate}

\noindent{\kaishu 证明.}

\noindent{\kaishu 存在性.}

假设有同态$p:A\to B$, 记$B'=\mathrm{Ran}\,p$, 那么:
\begin{enumerate}
    \item 对任意的常元$c$, $c^{B}=p(c^A)\in B'$,
    \item 对任意的$n$元函数符号$f$和$b_1,\cdots,b_n\in B'$,
    存在$a_1,\cdots,a_n$使得$p(a_1)=b_1,\cdots,p(a_n)=b_n$, 从而
    \[
        f^B(b_1,\cdots,b_n)=p(f^A(a_1,\cdots,a_n))\in B',
    \]
\end{enumerate}
于是由定理1.2有$B'\leqslant B$.
同时$p$可以看作$A$到$B'$的满射$q$, 所以$q:A\to B'$是满同态并且$p=\iota\circ q$.

\noindent{\kaishu 唯一性.} 显然. \hfill $\qedsymbol$

\vspace{20pt}





\section{子结构的语义}

\hypertarget{sec:定理3.1}{}
\noindent\textbf{定理3.1}

给定一阶语言$\mathcal{L}$和$\mathcal{L}$-结构$A,B$.
如果$A\leqslant B$, 则任给$A$上的变元赋值$\sigma$有:
\begin{enumerate}
    \item 对任意的项$t$, $t^{(A,\sigma)}=t^{(B,\sigma)}$,
    \item 对任意的无量词公式$\varphi$,
    $(A,\sigma)\models\varphi\leftrightarrow(B,\sigma)\models\varphi$.
\end{enumerate}

\noindent{\kaishu 证明.}

因为嵌入同态$\iota$是单同态, 由\href{02 结构的同态#nameddest=sec:定理3.3}
{单同态语义的定理}即得. \hfill $\qedsymbol$

\end{document}